\tituloI{PLANTEAMIENTO DEL ESTUDIO}

\tituloII{Planteamiento y formulación del problema}
\parrafoII{En la actualidad, el avance de la inteligencia artificial (IA) ha generado una transformación significativa en diversos ámbitos, especialmente en el sector educativo, donde su uso se ha expandido rápidamente entre estudiantes universitarios \cite{aparicio2023}.

Herramientas como asistentes conversacionales, generadores de texto y plataformas de apoyo a la programación han facilitado la realización de tareas académicas, el desarrollo de proyectos y la resolución de problemas complejos \cite{llanos2021}. Este fenómeno es particularmente relevante en carreras como Ingeniería de Sistemas e Informática, donde el uso de tecnologías digitales forma parte esencial del proceso de formación profesional, demostrando un alto índice de adopción tecnológica [1].

A nivel global, el uso de la inteligencia artificial en la educación superior ha crecido de manera exponencial, permitiendo optimizar el tiempo para que los estudiantes se concentren en procesos intelectuales de mayor relevancia, mejorando el acceso a la información y automatizando el aprendizaje \cite{sanchez2024}. Sin embargo, también ha generado preocupaciones relacionadas con la subordinación a la tecnología, la reducción del pensamiento crítico y la posible superficialidad en la adquisición del conocimiento autónomo \cite{chaparro2024}.

En el contexto peruano, si bien existe una creciente adopción de herramientas digitales en la educación superior, aún no se cuenta con lineamientos claros que garanticen un uso ético y confiable de la IA que maximice el potencial de los estudiantes \cite{unesco2024}.

En el ámbito local, específicamente en la sede Cusco de la Universidad Continental, se observa que los estudiantes de la carrera de Ingeniería de Sistemas e Informática utilizan cada vez más herramientas de inteligencia artificial para realizar sus actividades académicas, incluyendo tareas, proyectos e incluso programación \cite{llanos2021}. No obstante, no existe información sistematizada que permita comprender cómo, con qué frecuencia y con qué finalidad se emplean estas herramientas, ni cuál es la percepción de los estudiantes respecto a su utilidad y riesgos \cite{sanchez2024}.

Esta situación evidencia una problemática relacionada con la falta de conocimiento estructurado sobre el uso y la percepción de la inteligencia artificial, lo cual limita la toma de decisiones académicas informadas y la implementación de estrategias educativas adecuadas \cite{aparicio2023}. Asimismo, desde la perspectiva de la ingeniería de sistemas, esta problemática se vincula con la gestión de la información, el uso de tecnologías emergentes y la transformación digital de los procesos educativos \cite{llanos2021}.

Entre las posibles causas de esta problemática se encuentran la alta accesibilidad a estas herramientas, la ausencia de normativas institucionales claras sobre su uso académico \cite{unesco2024}, la necesidad de optimizar el tiempo por parte de los estudiantes y la falta de competencias digitales críticas para un uso adecuado \cite{chaparro2024}.

Como consecuencia, pueden generarse efectos como el aprendizaje superficial, la dependencia tecnológica, la disminución del desarrollo de habilidades cognitivas y la dificultad para evaluar el desempeño real de los estudiantes \cite{sanchez2024}. Si esta situación persiste, podría afectar la calidad de la formación profesional de los futuros ingenieros de sistemas, limitando su capacidad para resolver problemas de manera autónoma y crítica \cite{chaparro2024}.

En este contexto, se hace necesario desarrollar una investigación que permita analizar el uso y la percepción de la inteligencia artificial en los estudiantes, con el fin de generar información relevante que contribuya a un uso más adecuado y consciente de estas herramientas en el ámbito académico \cite{unesco2024}.}
\tituloIII{Problema General}
\parrafoIII{¿Cuál es el nivel de uso y la percepción de la inteligencia artificial en los estudiantes de la Escuela Académico Profesional de Ingeniería de Sistemas e Informática de la Universidad Continental, sede Cusco, 2026?}

\tituloIII{Problemas Específicos}
\begin{listaIII}
    \item ¿Cuáles son las herramientas de inteligencia artificial generativa más frecuentes utilizadas por los estudiantes en sus tareas de programación y desarrollo de software?
    \item ¿Cuál es la percepción de los estudiantes sobre la utilidad de la inteligencia artificial para el fortalecimiento de sus competencias profesionales?
    \item ¿Qué riesgos éticos y académicos perciben los estudiantes respecto al uso de asistentes virtuales en su formación académica?
\end{listaIII}


\tituloII{Objetivos}

\tituloIII{Objetivo General}
\parrafoIII{Determinar el nivel de uso y la percepción de la inteligencia artificial en los estudiantes de la Escuela Académico Profesional de Ingeniería de Sistemas e Informática de la Universidad Continental, sede Cusco, 2026.}

\tituloIII{Objetivos Específicos}
\begin{listaIII}
    \item Identificar las herramientas de inteligencia artificial generativa más frecuentes utilizadas por los estudiantes en sus tareas de programación y desarrollo de software.
    \item Analizar la percepción de los estudiantes sobre la utilidad de la inteligencia artificial para el fortalecimiento de sus competencias profesionales en ingeniería.
    \item Describir los riesgos éticos y académicos percibidos por los estudiantes respecto al uso de asistentes virtuales en su formación académica.    
\end{listaIII}

\tituloII{Justificación de la investigación}

\tituloIII{Justificación Teórica}
\parrafoIII{
    Desde el punto de vista teórico, esta investigación se justifica por la necesidad de generar nuevo conocimiento sobre la interacción entre el aprendizaje humano y los modelos de lenguaje de gran escala (LLM) en el contexto de la ingeniería. Según Vera \cite{vera2023}, la integración de la IA no es solo un cambio de herramienta, sino una transformación en el paradigma educativo que requiere ser documentada para comprender cómo se reconfiguran las competencias digitales de los futuros profesionales. El estudio aportará evidencia empírica que enriquecerá las teorías actuales sobre la adopción tecnológica y el aprendizaje autónomo en entornos digitales.
}

\tituloIII{Justificación Práctica}
\parrafoIII{
    En el plano práctico, la investigación es relevante porque proporcionará a la Universidad Continental, sede Cusco, un diagnóstico real sobre el uso de la IA en sus aulas. Esto es fundamental debido a que, a menudo, existe una brecha entre las herramientas que los estudiantes usan y las estrategias que los docentes implementan \cite{chaparro2024}. Los resultados permitirán a la facultad identificar patrones de riesgo, como la dependencia tecnológica, y oportunidades, como la optimización en el desarrollo de software, facilitando la creación de guías de buenas prácticas y lineamientos éticos internos.
}

\tituloIII{Justificación Metodológica}
\parrafoIII{
    Metodológicamente, el estudio se justifica mediante el diseño y validación de un instrumento de recolección de datos (cuestionario) adaptado específicamente a la realidad de los estudiantes de ingeniería en el Perú. Siguiendo a Hernández Sampieri et al. \cite{hernandez2014}, el valor de una investigación también radica en que sus procedimientos y herramientas puedan ser replicados o adaptados por otros investigadores interesados en medir el fenómeno de la inteligencia artificial en diferentes facultades o regiones, sirviendo como base para futuros estudios de mayor alcance.
}

\tituloII{Delimitación de la investigación}

\tituloIII{Delimitación Espacial}
\parrafoIII{
    El presente estudio se desarrollará físicamente en las instalaciones de la Universidad Continental, sede Cusco, ubicada en la \textbf{Vía de Evitamiento N° 801, distrito de San Jerónimo, Cusco}. La investigación se circunscribe exclusivamente a este campus universitario por ser el centro de formación de la población objetivo.
}

\tituloIII{Delimitación Temporal}
\parrafoIII{
    La fase de recolección de datos y análisis de resultados de la investigación se llevará a cabo en un periodo comprendido entre los meses de \textbf{julio y septiembre del año 2026}. Este intervalo permite capturar la actividad académica durante el ciclo correspondiente, asegurando la disponibilidad de los estudiantes para la aplicación de los instrumentos.
}

\tituloIII{Delimitación Social}
\parrafoIII{
    La investigación se centra específicamente en los estudiantes matriculados en la Escuela Académico Profesional de Ingeniería de Sistemas e Informática. Se considerará a alumnos desde el primer hasta el décimo ciclo, dado que el uso de herramientas de inteligencia artificial varía significativamente según el nivel de avance en la carrera y las competencias tecnológicas adquiridas \cite{continental2024}. No se incluirán en esta muestra a personal administrativo o docentes, con el fin de mantener la homogeneidad de la percepción estudiantil.
}

\tituloII{1.6 Hipótesis y Variables}

\tituloIII{Hipótesis}
\parrafoIII{
    De acuerdo con la naturaleza de la presente investigación, la cual posee un alcance \textbf{descriptivo}, no se ha formulado una hipótesis de trabajo. Según los lineamientos metodológicos para estudios que buscan especificar propiedades y características de un fenómeno sin establecer relaciones de causalidad, la formulación de hipótesis no es un requisito obligatorio.
}

\tituloIII{Variables}
\parrafoIII{
    En esta investigación se identifica una variable principal de estudio:
}

\begin{listaIII}
    \item \textbf{Variable única:} Uso y percepción de la inteligencia artificial.
\end{listaIII}

\tituloIII{Definición Operacional de la Variable}
\parrafoIII{
    La variable será medida a través de las siguientes dimensiones e indicadores:
}

\begin{listaIII}
    \item \textbf{Dimensión 1: Uso de la IA.} Indicadores: Frecuencia de uso, tipos de herramientas utilizadas (generadores de código, asistentes de texto), y propósito del uso (académico o personal).
    \item \textbf{Dimensión 2: Percepción de la IA.} Indicadores: Nivel de utilidad percibida, beneficios en el aprendizaje, y percepción de riesgos éticos/académicos.
\end{listaIII}